% Cote 119: offprints, reports, courses, CV and letters (typed and printed part): archive cote 119, site page 313 (PDF page 314).
% typescript. Draft compared: none.
% AI-transcribed by Claude Opus 5.5 (claude-opus-5-5, adaptive thinking, effort high), Message Batches API, batch r119_full (msgbatch_01DwZNm5Y7YsKBWyAN4sFs8w), 2026-11-13; not yet checked by eye.
\scanpage{119/313}
\mspage{-4-}

\uline{Géométrie algébrique}

[12] La théorie des classes de Chern, Bull. SMF 86
(1958), p. 137-154.

[13] Sur une note de Mattuck-Tate, J. Crelle 200
(1958), p. 137-154.

[14] The cohomology theory of abstract algebraic
varieties, Proc. Int. Congress, Edinburgh
(1958), p. 103-118.

[15] Elements de Géométrie Algébrique (rédigés avec
la coll. de Jean DIEUDONNE), Chap. I-IV, publ.
Math. IHES (1960/67), env. 1 800 pages.

[16] Séminaires de Géométrie Algébrique (SGA 1, ...,
7), IHES, 1960/69, env. 4 000 pages (en cours
de réédition chez Springer, Lecture Notes) :
SGA 1 Théorie du Groupe Fondamental
SGA 2 Cohomologie locale et Théorèmes de Lefschetz
locaux et globaux
SGA 3 (en coll. avec M. Demazure) Schémas en Groupes
des Topos et Cohomologie étale des Schémas
SGA 5 Cohomologie l-adique et fonctions L
SGA 6 (en coll. avec J. Berthelot et J.L. Illusie)
Théorie des Intersections et Théorèmes de
Riemann-Roch
SGA 7 Groupe de Monodromie en Géométrie Algébrique

[17] Un théorème sur les homomorphismes de schémas
abéliens, Invent. Math. 2 (1966), p. 59-78.

[18] Dix exposés sur la cohomologie des schémas
(en coll. avec J. Giraud, S. Kleiman,
M. Raynaud, J. Tate), North Holland, 1968.

[19] Catégories cofibrées additives et complexe
cotangent relatif, Lecture Notes in Maths.,
Springer n° 79 (1968), 167 pages.

\hand{[20] Platitude \unsure{d'une|d'un} adhérence schématique}
\hand{et lemme de Hironaka \unsure{généralisé} (avec}
\hand{H. \unsure{Seydi}), manuscripta math., 5, 323-339,}
\hand{(1971)}

\margin{318}
